% !TEX TS-program = lualatex
% OSCAR organigram.  Redraw of the original, a 2020 Keynote drawing that
% was later edited in Affinity Designer.  Build with `make`, or:
%
%     lualatex OSCAR-organigram.tex
%
% Coordinates are in millimetres, y pointing up.  Everything you are
% likely to change lives in the two configuration blocks below.
\documentclass[border=3mm]{standalone}

% Byte-reproducible output: no timestamps, and a fixed trailer ID in place of
% the random one, so the committed PDF changes only when the drawing does.
\pdfvariable suppressoptionalinfo=127
\pdfvariable trailerid{[<5be15cfc1fdee8c6e98f233ff4586422> <5be15cfc1fdee8c6e98f233ff4586422>]}

\usepackage{fontspec}
% Medium as the upright weight, as in the original (which used Helvetica Neue).
\setsansfont{FiraSans}[Extension=.otf, UprightFont=*-Medium, BoldFont=*-Bold]
\renewcommand{\familydefault}{\sfdefault}

\usepackage{tikz}
\usetikzlibrary{arrows.meta,calc,decorations.pathreplacing}

% ---------------------------------------------------------------- colours
% Blue from the website (#069), amber and slate from the OSCAR logo.
\definecolor{topicblue} {HTML}{0A5F86}   % the five mathematical areas
\definecolor{juliablue} {HTML}{AFD5E8}   % Oscar.jl backdrop
\definecolor{wrapgreen} {HTML}{B7DE96}   % Julia wrapper packages
\definecolor{coregreen} {HTML}{38663C}   % cornerstone systems
\definecolor{subgreen}  {HTML}{55984F}   % components inside a cornerstone
\definecolor{basegrey}  {HTML}{D9D9DC}   % third-party libraries
\definecolor{baseamber} {HTML}{FAB916}   % third-party libraries co-developed by us
\definecolor{arrowcol}  {HTML}{B4622D}   % dependency arrows
\definecolor{framecol}  {HTML}{1A1A1A}   % the OSCAR frame

% ------------------------------------------------------------- dimensions
\def\ncols{5}            % number of wrapper/cornerstone columns
\def\ntopics{5}          % number of mathematical areas
\def\colw{52.9}          % column width
\def\colgap{1.8}
\def\topicw{45}          % width of an area box
\def\topicgap{8}
\def\basegap{4.5}        % gap between the boxes of the bottom row

\def\yOscarTop{158}      % Oscar.jl backdrop
\def\yOscarBot{88.5}
\def\yTopic{144.5}       % top edge of the area row
\def\hTopic{42.5}
\def\yWrap{86.5}         % top edge of the wrapper row
\def\hWrap{12.7}
\def\yCorner{66.5}       % top edge of the cornerstone row
\def\hCorner{42.5}
\def\yBase{4}            % top edge of the foundations row
\def\hBase{14}
\def\yFrameTop{173.5}    % OSCAR frame
\def\yFrameBot{20}
\def\yDashJulia{70.5}    % Julia / cornerstone divider
\def\subgap{2.9}         % between components inside a cornerstone
\def\titleh{12}          % space a cornerstone title takes at the top
\def\rad{3.5}            % corner radius, boxes
\def\radbig{5}           % corner radius, the two containers

\pgfmathsetmacro{\stackw}{\ncols*\colw+(\ncols-1)*\colgap}
\pgfmathsetmacro{\hStack}{\yWrap-\yCorner+\hCorner}    % wrapper row + cornerstone row
\pgfmathsetmacro{\topicmargin}{(\stackw-\ntopics*\topicw-(\ntopics-1)*\topicgap)/2}
\newcommand{\colx}[1]{(#1-1)*(\colw+\colgap)}          % left edge of column #1
\newcommand{\topicx}[1]{\topicmargin+(#1-1)*(\topicw+\topicgap)}

% ----------------------------------------------------------------- styles
\tikzset{
  box/.style     = {anchor=north west, rounded corners=\rad mm, align=center,
                    inner sep=0pt, text=white},
  topic/.style   = {box, fill=topicblue, minimum width=\topicw mm,
                    minimum height=\hTopic mm, font=\fontsize{16}{19}\selectfont},
  wrap/.style    = {box, fill=wrapgreen, text=black, minimum width=\colw mm,
                    minimum height=\hWrap mm, font=\fontsize{20}{20}\selectfont},
  core/.style    = {box, fill=coregreen, minimum width=\colw mm,
                    minimum height=\hCorner mm},
  coretitle/.style = {anchor=north, yshift=-3mm, inner sep=0pt, text=white,
                    font=\fontsize{20}{20}\selectfont},
  sub/.style     = {box, fill=subgreen, anchor=north,
                    minimum width=40mm, minimum height=9.8mm,
                    font=\fontsize{14}{14}\selectfont},
  listbox/.style = {box, fill=subgreen, anchor=north, align=left,
                    text width=32mm, inner xsep=4mm, inner ysep=2mm},
  base/.style    = {box, fill=basegrey, text=black,
                    minimum width=38mm, minimum height=\hBase mm,
                    font=\fontsize{20}{20}\selectfont},
  ours/.style    = {base, fill=baseamber},
  frame/.style   = {rounded corners=\radbig mm, line width=0.7mm, draw=framecol},
  backdrop/.style= {rounded corners=\radbig mm},
  grouplabel/.style = {anchor=north west, inner sep=0pt, font=\fontsize{20}{20}\selectfont},
  sidelabel/.style  = {anchor=east, inner sep=0pt, font=\fontsize{24}{24}\bfseries},
  divider/.style = {dashed, dash pattern=on 2.1mm off 2.1mm},
  brace/.style   = {decorate, decoration={brace, amplitude=5mm},
                    line width=0.6mm, line cap=round, line join=round},
  dep/.style     = {arrowcol, line width=0.8mm, line cap=butt,
                    -{Triangle[length=2.8mm,width=2.8mm]}},
}

% \dropto{<from>}{<to>}          straight down
% \routeto{<from>}{<depth>}{<to>}  down by <depth>, sideways, then down
\newcommand{\dropto}[2]{\draw[dep] (#1) -- (#1 |- #2);}
\newcommand{\routeto}[3]{\draw[dep] (#1) -- ++(0,{-(#2)}) -| (#3);}

% \complist{<heading>}{<items, separated by \\>} -- contents of a listbox
\newcommand{\complist}[2]{\fontsize{12}{14}\selectfont\bfseries #1\par
  \fontsize{10.5}{12.5}\selectfont\mdseries #2}
\newcommand{\li}{\textbullet~}

\begin{document}
\begin{tikzpicture}[x=1mm, y=1mm]

% ============================================================ the OSCAR box
\draw[frame] (-3.5,\yFrameTop) rectangle (\stackw+3.5,\yFrameBot);
\node[grouplabel] at (3,\yFrameTop-4) {OSCAR};

\fill[backdrop, juliablue] (0,\yOscarTop) rectangle (\stackw,\yOscarBot);
\node[grouplabel] at (3,\yOscarTop-4) {Oscar.jl};

% ================================================== mathematical areas
\node[topic] (nt)   at ({\topicx{1}},\yTopic) {number theory};
\node[topic] (grp)  at ({\topicx{2}},\yTopic) {group and\\representation\\theory};
\node[topic] (trop) at ({\topicx{3}},\yTopic) {tropical and\\polyhedral\\geometry};
\node[topic] (ag)   at ({\topicx{4}},\yTopic) {algebraic\\geometry and\\commutative\\algebra};
\node[topic] (nc)   at ({\topicx{5}},\yTopic) {non-\\commutative\\algebra and\\free probability\\theory};

% ============================================== Julia wrapper packages
% Column 1 is pure Julia all the way down, hence one tall box.
\node[wrap, minimum height=\hStack mm,
      top color=wrapgreen, bottom color=coregreen] (aa) at ({\colx{1}},\yWrap) {};
\node[wrap] (gapjl)  at ({\colx{2}},\yWrap) {GAP.jl};
\node[wrap] (pmjl)   at ({\colx{3}},\yWrap) {Polymake.jl};
\node[wrap] (singjl) at ({\colx{4}},\yWrap) {Singular.jl};
\node[wrap, font=\fontsize{15}{15}\selectfont] (asjl) at ({\colx{5}},\yWrap) {AlgebraicSolving.jl};

% ================================================= cornerstone systems
\node[core] (gap)  at ({\colx{2}},\yCorner) {};  \node[coretitle] at (gap.north)  {GAP};
\node[core] (pm)   at ({\colx{3}},\yCorner) {};  \node[coretitle] at (pm.north)   {polymake};
\node[core] (sing) at ({\colx{4}},\yCorner) {};  \node[coretitle] at (sing.north) {Singular};

\node[sub] (hecke)    at (aa.north |- 0,\yCorner)              {Hecke.jl};
\node[sub] (nemo)     at ($(hecke.south)-(0,\subgap)$)         {Nemo.jl};
\node[sub, font=\fontsize{12}{12}\selectfont]
           (abstract) at ($(nemo.south)-(0,\subgap)$)          {AbstractAlgebra.jl};
\node[listbox] (gappkg) at ($(gap.north)-(0,\titleh)$)
      {\complist{packages}{\li atlasrep\\\li ctbllib\\\li transgrp\\\li\textellipsis}};
\node[listbox] (pmapp)  at ($(pm.north)-(0,\titleh)$)
      {\complist{applications}{\li fan\\\li polytope\\\li tropical\\\li\textellipsis}};

\node[sub] (plural)   at ($(sing.north)-(0,\titleh)$)          {PLURAL};
\node[sub] (letter)   at ($(plural.south)-(0,\subgap+1.7)$)    {Letterplace};

% ============================================== third-party foundations
\node[base] (gmp)   at (-1.4,\yBase) {GMP};
\node[ours] (flint) at ($(gmp.north east)+(\basegap,0)$)   {FLINT};
\node[base] (cdd)   at ($(flint.north east)+(\basegap,0)$) {cdd};
\node[base] (bliss) at ($(cdd.north east)+(\basegap,0)$)   {bliss};
\node[base, minimum width=45.5mm]
            (gpi)   at ($(bliss.north east)+(\basegap,0)$) {GPI-Space};
\node[ours] (msolve) at ($(gpi.north east)+(\basegap,0)$) {msolve};
\node[anchor=west, inner sep=0pt, font=\fontsize{20}{20}\selectfont]
      at ($(msolve.east)+(3,0)$)
      {\textperiodcentered\,\textperiodcentered\,\textperiodcentered};

% ========================================================= layer dividers
% The frame is the OSCAR / outside-world boundary, so no dashed line there.
\draw[divider] (-32,\yDashJulia) -- (\stackw+3,\yDashJulia);

\draw[brace] (-19,\yWrap-\hWrap) -- (-19,\yOscarTop);
\node[sidelabel] at (-26,{(\yWrap-\hWrap+\yOscarTop)/2}) {Julia};
\draw[brace] (-19,\yCorner-\hCorner) -- (-19,\yCorner);
\node[sidelabel] at (-26,{\yCorner-\hCorner/2})          {Cornerstones};
\draw[brace] (-19,\yBase-\hBase) -- (-19,\yBase);
\node[sidelabel] at (-26,{\yBase-\hBase/2})              {Foundations};

\node[anchor=east, inner sep=0pt, font=\fontsize{13}{13}\selectfont]
      (legend) at (\stackw+3.5,{\yBase-\hBase-6}) {co-developed within OSCAR};
\fill[baseamber, rounded corners=1mm]
      ([xshift=-3mm,yshift=-2.2mm]legend.west) rectangle ++(-7,4.4);

% ================================================== areas -> wrappers
\dropto {[xshift=-16.8mm]nt.south}{aa.north}
\routeto{[xshift=-10.8mm]nt.south}{9.4}{[xshift=-19.4mm]gapjl.north}
\routeto{[xshift=-4.8mm]nt.south}{7.1}{[xshift=-19.9mm]pmjl.north}

\routeto{[xshift=-9.1mm]grp.south}{1.4}{[xshift=14.2mm]aa.north}
\dropto {[xshift=-0.9mm]grp.south}{gapjl.north}

\dropto {[xshift=1.3mm]trop.south}{pmjl.north}

\routeto{[xshift=-19mm]ag.south}{3.9}{[xshift=19.1mm]aa.north}
\routeto{[xshift=-6.4mm]ag.south}{6}{[xshift=20.2mm]pmjl.north}
\dropto {[xshift=7.8mm]ag.south}{singjl.north}
\routeto{[xshift=14.5mm]ag.south}{4.7}{[xshift=4.5mm]asjl.north}

\routeto{[xshift=-2mm]nc.south}{7.8}{[xshift=13.8mm]singjl.north}

% ============================================ wrappers -> cornerstones
\dropto{gapjl.south}{gap.north}
\dropto{pmjl.south}{pm.north}
\dropto{singjl.south}{sing.north}
\dropto{asjl.south}{msolve.north}

% ======================================= cornerstones -> foundations
\dropto {[xshift=20.55mm]aa.south}{flint.north}
\routeto{[xshift=-22.85mm]pm.south}{8}{[xshift=7.9mm]flint.north}
\routeto{[xshift=-16.85mm]pm.south}{12}{[xshift=7.4mm]cdd.north}
\dropto {[xshift=1.15mm]pm.south}{bliss.north}
\routeto{[xshift=-23.55mm]sing.south}{16}{[xshift=13.9mm]flint.north}
\dropto {[xshift=-5.55mm]sing.south}{gpi.north}

\end{tikzpicture}
\end{document}
